\chapter{引言}

\section{研究背景}
教育数字化背景下，编程教育逐渐成为中小学生信息素养培养的重要内容。编程学习不仅涉及语法掌握，还关系到逻辑推理、抽象建模和问题求解能力的形成。然而，初学者在语法理解、报错定位和任务分解方面容易受阻，教师也难以在有限课堂时间内为每名学生提供持续反馈。

大语言模型在自然语言理解、代码解释和交互式问答方面的能力，为编程学习平台提供了新的支持方式。已有研究表明，LLM 可用于代码生成、调试辅助、形成性反馈和个性化学习支持。但在中小学场景中，模型回答的准确性、适龄性和可控性仍需谨慎处理，不能简单把通用问答能力直接搬入课堂。

因此，本文关注的不是单一大模型接口调用，而是面向中小学生编程学习流程，设计并实现一套集课程学习、在线编程、智能辅学、教师干预和过程记录于一体的平台。平台以“教学约束、风险兜底、过程留痕”为设计前提，使智能能力服务于真实学习过程。

\section{研究意义}
本文研究意义主要体现在两方面：一是现有 LLM 赋能编程教育研究多聚焦高等教育或单一任务，对面向中小学生、覆盖课程学习到过程记录的平台化研究相对不足；二是本文将大模型能力嵌入日常学习流程，而非作为孤立问答工具使用，强调教师可控、过程可追踪和输出可约束。

\section{国内外研究现状}
国内外相关研究主要集中在代码生成、代码解释、错误诊断、形成性反馈和学习行为分析等方向。已有综述表明，LLM 在学生支持、教师辅助和自适应学习方面具有潜力，但也伴随可靠性、偏差与伦理风险。面向 K-12 的研究同样提示：教学场景需要教师在环、过程监督和结果复核，平台设计不应退化为单一聊天窗口。

\section{研究内容与目标}
本文围绕以下目标展开：构建学生与教师双角色编程学习平台；实现课程、知识点、任务、资源和行为记录等业务模块；集成大模型能力并实现代码解释、报错分析、问题分解、抽象建模与算法设计引导；实现在线代码运行与安全控制；完成功能与运行效果测试分析。

\section{论文结构安排}
全文共七章：第一章为引言；第二章为概述；第三章为系统总体设计；第四章为系统详细设计；第五章为软件实现；第六章为系统测试与调试；第七章为结论。

\chapter{概述}

\section{系统建设目标概述}
系统面向中小学生编程学习与教师教学管理两个侧面。学生侧聚焦课程学习、在线练习、报错理解和智能提示；教师侧聚焦资源维护、过程观察、消息互动和教学干预。平台将大模型能力嵌入学习关键环节，而不是作为独立功能孤岛。

\section{大语言模型与智能辅学}
LLM 在教育场景可承担问答辅助、内容解释、错误分析和学习引导等任务，但并不等同于替代教师。本文将 LLM 定位为过程性辅学模块，强调在学生卡点处提供适量提示、在错误定位处提供结构化解释，并保留教师审阅、追踪和干预空间。

\section{前端开发技术}
前端采用 Nuxt 3 + Vue 3 组合式 API，配合 Tailwind CSS、DaisyUI 与 TypeScript。该组合兼顾交互开发效率、组件化组织能力与工程可维护性。

\section{后端开发技术}
后端采用 FastAPI，数据库访问通过 SQLAlchemy 实现 ORM 映射，认证采用 JWT，缓存采用 Redis。整体后端结构强调接口分层、权限边界和可扩展性。

\section{数据存储与资源管理技术}
结构化数据存储在 SQL Server，包括用户、课程、任务、知识点、资源、评论、点赞与学习记录等。非结构化资源通过阿里云 OSS 存储，提升教学资源访问与管理效率。

\section{在线编程与安全执行技术}
平台实现在线代码运行，执行前进行语法检查和安全检查，过滤危险导入与危险函数调用，并通过超时控制和执行次数限制降低恶意代码与死循环风险。

\section{系统需求分析}
\subsection{应用场景分析}
平台主要覆盖课前准备、课中讲解、课后巩固和过程评价等教学环节。教师可配置任务与资源，学生可完成代码编写、运行调试和即时提问，教师可依据过程数据进行后续教学调整。

\subsection{用户角色分析}
系统核心角色包括学生与教师。学生侧关注流程清晰、反馈及时和解释易懂；教师侧关注组织能力、统计能力和教学干预能力。权限设计上，学生以学习与提交为主，教师拥有课程维护、资源管理与记录查看等更高权限。

\subsection{功能需求分析}
学生端需求包括登录认证、课程任务查看、在线编程运行、智能问答与错误解析、消息互动、学习记录查看。教师端需求包括进度查看、错误统计、资源管理、课程任务维护与消息交流。智能辅学需覆盖自由问答、代码解释、报错分析、问题分解、抽象建模与算法设计引导。

\subsection{非功能需求分析}
平台需满足性能、安全、可扩展与易用性要求：在课堂高峰场景保持可接受响应；保障认证与代码执行安全；为模型替换、知识库增强与个性化推荐预留接口；界面与反馈贴近中小学生认知水平。

\subsection{可行性分析}
从技术、经济和运行三方面均具备可行性：技术栈成熟，开发与部署成本可控，且与学校常见教学流程具有较高匹配度。

\chapter{系统总体设计}

\section{系统设计目标}
总体设计遵循模块化、可扩展与安全可控原则，覆盖课程学习、在线编程、师生互动等基础能力，并通过大模型增强学习反馈的及时性与针对性。

\section{系统总体架构设计}
系统采用前后端分离架构：前端负责页面展示与交互；后端承担业务处理、认证鉴权、代码执行、数据持久化与智能能力调度。前端通过统一接口访问后端服务，实现教学流程闭环。

\section{系统功能模块设计}
平台按业务域划分模块：用户管理、课程管理、任务管理、消息互动、资源管理、行为记录与智能辅学。模块边界清晰，便于独立维护与后续扩展。

\section{前后端结构设计}
\subsection{前端结构设计}
前端采用“页面 + 组件”组织方式。页面层承载业务流程，组件层沉淀可复用交互单元。学生端强调学习主路径简洁，教师端强调内容维护效率与信息聚合能力。

\subsection{后端结构设计}
后端按业务域组织路由，公开接口与鉴权接口分组管理。认证、异常处理和统一返回机制尽量集中沉淀，业务逻辑按模块解耦，兼顾一致性与扩展弹性。

\chapter{系统详细设计}

\section{数据库设计}
\subsection{数据库概念结构设计}
数据库围绕用户、课程、任务、知识点、资源、评论、点赞与学习记录等核心实体展开，既保存结果数据，也保留过程数据，为统计分析与教学干预提供基础。

\subsection{主要数据表设计}
主要数据表包括 `users`、`courses`、`contents`、`tasks`、`knowledges`、`resources`、`comments`、`likes`、`records`。各表通过主外键或业务关联字段支撑课程学习、任务执行和行为分析。

\section{智能辅学模块设计}
智能辅学链路包括“提问输入—场景识别—提示词构造—模型生成—结果返回”。为避免偏离教学目标，模块进一步细化为“场景识别—约束生成—教学化输出—风险兜底”四层结构，强调适龄表达、分步引导和高风险输出兜底。

\section{系统部署设计}
前端通过 Node 环境部署 Nuxt，后端通过 Uvicorn 运行 FastAPI；数据库采用 SQL Server，缓存采用 Redis，非结构化资源存储于阿里云 OSS。该部署方式职责划分清晰，满足毕业设计规模下的联调与演示需求。

\chapter{软件实现}

\section{用户认证与权限控制实现}
系统采用 JWT 认证机制，前端执行页面级拦截，后端执行接口级鉴权与角色校验。该模块保证“身份可信、访问有边界”，是平台稳定运行基础。

\section{课程与知识点管理功能实现}
平台通过课程分类选择、课程搜索与知识点选择器组织内容。教师可完成课程编辑、知识点关联与内容维护，后端保证结构化数据持久化与关系一致性。

\section{在线编程与代码运行功能实现}
代码运行流程为输入接收、静态检查、受控执行与结果回传。系统兼顾执行边界控制与教学反馈可用性，并将运行过程写入行为记录，为教师分析学习状态提供依据。

\section{大模型辅助教学功能实现}
\subsection{自由问答功能实现}
自由问答支持编程相关流式问答，通过系统提示词约束回答范围与语气，避免偏离课程目标。

\subsection{代码解释功能实现}
代码解释按“整体功能—关键逻辑—局部语句”组织输出，帮助初学者形成由整体到局部的理解路径。

\subsection{报错解析功能实现}
报错解析采用“三步式”结构：错误含义、错误定位、修改引导，降低初学者理解门槛并提升可操作性。

\subsection{问题分解功能实现}
系统将综合任务拆解为可执行步骤，帮助学生从“无从下手”转为“分步推进”。

\subsection{抽象建模功能实现}
模块引导学生从具体问题提炼一般规律，逐步形成可迁移的求解框架。

\subsection{算法设计引导功能实现}
系统通过固定流程模板（开始、输入、处理、循环、输出、结束）帮助学生形成结构化算法表达能力。

\section{师生互动与消息功能实现}
平台支持消息发送、读取、未读状态管理与对话展示，用于课后答疑与学习交流，补足“人际反馈链路”。

\section{教学资源管理功能实现}
教师端支持资源创建、修改、删除与权限设置，并与课程和知识点建立关联，提升资源复用效率。

\section{学习行为记录与事件追踪功能实现}
系统记录页面访问、代码运行与错误类型等数据，并通过可视化页面呈现，为教学干预与后续研究提供证据基础。

\chapter{系统测试与调试}

\section{测试环境}
测试环境覆盖 Nuxt 3 + Node.js 前端、FastAPI + Python 后端、SQL Server、Redis 与模型接口环境，尽量贴近真实运行条件。

\section{功能测试}
功能测试围绕登录认证、在线编程、报错分析、消息互动、资源管理等核心链路开展，除验证单点可用性外，还检查流程闭环与结果回传一致性。

\section{性能测试}
性能测试关注接口响应时间、代码运行耗时、模型响应时延与并发访问表现。当前阶段以系统设计与实现为主，性能测试作为后续补充验证项。

\section{测试结果分析}
当前联调结果显示核心功能链路已可闭环运行，架构设计与模块划分具备可行性。后续将结合真实教学实验与量表指标，对学习效果与教学支持效果进行更系统验证。

\chapter{结论}

本文围绕中小学生编程学习场景，完成了基于大模型的智能编程学习平台设计与实现。系统覆盖课程管理、任务学习、在线编程、智能辅学、消息互动、资源管理与行为记录等模块，并通过场景划分、提示约束与输出组织实现教学化落地。

与直接接入聊天接口的实现方式相比，本文更强调教学适配、系统约束与教师可控性。平台重点不是替代教师，而是构建教师、学生与大模型协同参与的学习支持机制。

现阶段结论主要集中在系统设计合理性与核心功能实现层面。平台对学习兴趣、自我效能感和长期学习行为的影响仍需通过真实课堂实验、对照测试与过程数据分析进一步验证。后续可在知识库增强、学习者建模、提示策略评估和教师干预机制等方向继续拓展。
